\documentclass{article}

\usepackage[final]{neurips_2019}

\usepackage[utf8]{inputenc}
\usepackage[T1]{fontenc}
\usepackage{hyperref}
\usepackage{url}
\usepackage{booktabs}
\usepackage{amsfonts}
\usepackage{nicefrac}
\usepackage{microtype}
\usepackage{graphicx}
\usepackage{xcolor}
\usepackage{lipsum}
\usepackage{amsmath}
\usepackage{graphicx}   % For images
\usepackage{float}      % For [H] positioning
\usepackage{booktabs}   % For table lines (toprule, midrule)
\usepackage{caption}    % For caption formatting
% \usepackage{subcaption} % (Optional) If you revert to using 'subfigure' instead of minipage
\newcommand{\note}[1]{\textcolor{blue}{{#1}}}

\title{
  Title of your project \\
  \vspace{1em}
  \small{\normalfont Columbia COMS4705 Final Project Report} \\
  \small{\normalfont \textbf{Keywords:} \textit{keyword1, keyword2, ...}}
}

\author{
  Joshua Hegstad \\
  Department of Computer Science \\
  Columbia University \\
  \texttt{jlh2288@columbia.edu} \\
  \And
  Ahmed Jaber \\
  Department of Computer Science \\
  Columbia University \\
  \texttt{amj2234@columbia.edu} \\
  \And
  Farhaan Siddiqui \\
  Department of Computer Science \\
  Columbia University \\
  \texttt{fs2872@columbia.edu} \\
}

\begin{document}

\maketitle

% \begin{center}
%     \note{This template is built on NeurIPS 2019 template\footnote{\url{https://www.overleaf.com/latex/templates/neurips-2019/tprktwxmqmgk}} and adapted from Stanford CS224N Natural Language Processing with Deep Learning 
%     }
% \end{center}

\begin{abstract}
Closed-book question answering (QA) systems must answer questions using only parametric knowledge, which makes them prone to \emph{closed-book hallucination}: confident but factually incorrect answers with no supporting evidence \cite{ji2023survey}. We view this behavior as arising from two failure modes: (1) the model lacks the relevant knowledge, and (2) the knowledge is implicitly present in its parameters, but the model does not reliably retrieve or apply it at inference time \cite{kadavath2022language,lin2022teaching}. We study these failure modes on HotPotQA using Qwen2.5-Instruct-7B and propose a two-part mitigation pipeline.

To address the ``no knowledge'' case, we train the model to abstain instead of hallucinating. We run a chain-of-thought (CoT) teacher on the training split, replace each incorrect prediction with the phrase ``I don't know,'' and finetune Qwen2.5-7B on an 80/20 mix of correct answers and abstentions. We then apply verifier-based reward learning that combines an NLI-derived factuality score with an entropy-based confidence term, encouraging calibrated abstention and discouraging confident errors. To address the ``cannot access knowledge'' case, we perform CoT distillation from a larger Qwen3-235B teacher, following recent work showing that CoT distillation can transfer reasoning behavior to smaller models \cite{li2023cotdistill}. On HotPotQA, this pipeline reduces hallucinated responses, increases appropriate ``I don't know'' answers, and yields modest EM/F1 gains over a supervised-only baseline, suggesting that jointly training when to answer and when to abstain improves closed-book reliability.
\end{abstract}




\section{Key Information to include}
\begin{itemize}
    \item Mentor: Melody Ma
    \item External Collaborators (if you have any): None
    \item Sharing project: N/A
\end{itemize}

% {\color{red} This template does not contain the full instruction set for this assignment; please refer back to the milestone instructions PDF.}

\section{Introduction}
\section{Introduction}
Large Language Models (LLMs) have demonstrated remarkable capabilities in knowledge retrieval and reasoning. However, when operating in "closed-book" settings—where they must rely solely on parametric knowledge—they frequently suffer from \emph{hallucination}: generating plausible-sounding but factually incorrect assertions \cite{ji2023survey}. This issue is particularly acute in Small Language Models (SLMs) (e.g., < 10B parameters), which have limited capacity to store world knowledge and often lack the calibration to distinguish between what they know and what they have forgotten or never learned. 

This unreliability poses a critical barrier to deployment in safety-critical domains where admitting ignorance is preferable to confident error. While Retrieval-Augmented Generation (RAG) mitigates this by grounding generation in external documents, many applications (e.g., edge devices, privacy-constrained environments) necessitate reliable closed-book performance. Thus, a key open challenge is: \emph{How can we teach small models to reason effectively about what they know, and robustly abstain when they do not?}

We address this challenge via a three-stage pipeline designed to enhance both capability and calibration in Qwen2.5-7B-Instruct. First, to mitigate the "cannot access knowledge" failure mode, we employ **Chain-of-Thought (CoT) Distillation** from a massive teacher model (Qwen3-235B). We hypothesize that explicit reasoning traces help the model "unlock" latent parametric knowledge that is inaccessible via direct prompting. Second, to address the "no knowledge" failure mode, we implement **Abstention-Aware Fine-Tuning**, teaching the model to explicitly output "I don't know" for questions that induce hallucination. Finally, we refine this behavior using **Reinforcement Learning from Verifier Feedback (RLVF)**. We construct a reward function that balances factuality (verified by an NLI model) with confidence (measured via entropy), incentivizing the model to be confident only when correct.

Our main contributions are as follows:
\begin{enumerate}
    \item We demonstrate that CoT distillation significantly outperforms standard SFT in closed-book QA, improving Exact Match (EM) accuracy from 18.86\% to 24.31\% on HotPotQA.
    \item We introduce a supervised abstention baseline that achieves near-perfect precision (97.5\%) in detecting its own knowledge boundaries, though at the cost of recall.
    \item We show that RLVF with asymmetric penalties effectively tunes the risk-coverage trade-off, increasing answer rate by $\sim5\%$ while maintaining safety.
    \item We validate that unsupervised uncertainty metrics—specifically \emph{Semantic Entropy (SE)}—serves as a robust proxy for correctness, allowing for efficient rejection of remaining hallucinations at inference time.
\end{enumerate}

\section{Related Work}
\section{Related Work}

\textbf{Hallucination in LLMs}: Hallucination refers to the generation of text that is nonsensical or unfaithful to the provided source or established world knowledge \cite{ji2023survey}. In closed-book settings, where models rely on parametric memory, hallucinations often stem from the model's inability to accurately assess the boundaries of its own knowledge \cite{kadavath2022language}. While larger models exhibit emergent self-calibration, Small Language Models (SLMs) are particularly prone to confident errors, necessitating explicit intervention strategies.

\textbf{Reasoning Distillation}: Chain-of-Thought (CoT) prompting significantly enhances LLM reasoning but typically requires large model scales (>100B parameters) to be effective. To transfer this capability to smaller models, \cite{li2023cotdistill} and \cite{ho2023largelanguagemodelsreasoning} proposed CoT distillation, where a student model is fine-tuned on reasoning traces generated by a large teacher. Our work extends this by imposing "internalized knowledge" constraints on the teacher, ensuring the distilled reasoning relies on facts the student can plausibly retrieve from its own weights rather than external context.

\textbf{Abstention and Uncertainty}: Teaching models to say "I don't know" is a classic problem in reliable QA \cite{rajpurkar2018know}. Recent approaches have focused on unsupervised methods to detect uncertainty, such as analyzing the entropy of semantic clusters in model generations (Semantic Entropy) \cite{Farquhar2024Detecting} or inspecting internal attention patterns (RAUQ) \cite{vazhentsev2025rauq}. \cite{lin2022teaching} demonstrated that models can be trained to verbalize this uncertainty. We combine these streams by using supervised abstention for base calibration and refining it with NLI-based Reinforcement Learning (RLVF) that directly rewards the alignment between confidence (low entropy) and factuality.

\section{Approach}
This section details your approach(es) to the problem. 
For example, this is where you describe the architecture of your neural network(s), and any other key methods or algorithms.

\textbf{Key Models:}
\begin{itemize}
\item Qwen2.5-7B-Instruct \cite{qwen2025qwen25technicalreport}: Referred to as\textbf{ Qwen2.5-Instruct}.
\item Qwen2.5-7B-Instruct (SFT): Referred to as \textbf{Qwen2.5-SFT}. Finetuned using LoRA on 10,000 question-answer pairs in HotPotQA 
\item Qwen2.5-7B-Instruct (Finetune-CoT): Referred to as \textbf{Qwen2.5-FCoT}. Finetuned on Dataset B to generate a reasoning trace and answer for each question, without the context of supporting paragraphs. 
\item Qwen2.5-7B-Instruct (Abstention Finetuned): Referred to as \textbf{Qwen2.5-Abstain}. Finetuned using LoRA on the question-answer/abstention pairs in Dataset C. The dataset is imbalanced to 80\% correct answers and 20\% abstentions to prevent the model from collapsing to constant abstention.
\end{itemize}

\subsection{Finetune-CoT} \label{sec:finetune-cot}
To enable complex reasoning capabilities in our smaller student model, we employed Fine-tune-CoT \cite{ho2023largelanguagemodelsreasoning}, a framework that leverages the capabilities of large language models (LLMs) to supervise the training of smaller models. While massive LLMs can solve complex tasks using Chain-of-Thought (CoT) prompting, they are often computationally prohibitive to deploy. Fine-tune-CoT circumvents this by using a large "Teacher" model to generate synthetic reasoning traces, which are then used to fine-tune a "Student" model. This effectively transfers the reasoning ability, allowing the smaller model to learn how to solve the problem.

The standard Fine-tune-CoT pipeline consists of three stages: (1) prompting the teacher to generate step-by-step rationales, (2) curating the data by filtering out samples with incorrect answers, and (3) fine-tuning the student on the validated reasoning traces.

In our implementation, we adapted the Generation and Curation stages to create a specialized "Internalized Knowledge" dataset (Dataset B). We utilized Qwen3-235B as our teacher model. To address the common issue of hallucination during reasoning generation, we provided the teacher with the ground-truth paragraphs as context. However, to ensure the student model learns to function in a closed-book setting, we imposed strict constraints on the teacher's output: it was prompted to frame its reasoning as a step-by-step recall of facts, explicitly forbidding citation markers (e.g., "according to the passage"). We then implemented a validation loop that filtered samples not only for answer correctness (as in standard Fine-tune-CoT ) but also for adherence to this "internalized" stylistic constraint.

\subsection{Abstention-Aware Fine-Tuning} \label{sec:abstension}
\textbf{Abstention-Aware Fine-Tuning: }To mitigate hallucinations, we employed a specialized fine-tuning strategy designed to calibrate the model's confidence. Rather than forcing the model to answer every query—which encourages hallucination on out-of-distribution or forgotten knowledge—we fine-tuned the model to explicitly output a refusal phrase ("I don't know") when it lacks the internal knowledge to answer correctly. 

We first aimed to map the "knowledge boundary" of our baseline model (Qwen2.5-FCoT). We hypothesize that if the model generates an incorrect answer for a training example $x$ despite having access to its internal weights, that example represents a "known unknown" or a hallucination-prone region. We performed inference on the first $N=40,000$ examples of the HotpotQA training set. For every example, we compared the model's predicted answer $\hat{y}$ against the ground truth $y$ using Exact Match (EM) scoring. This effectively partitioned our training data into two sets: 
\begin{itemize}
    \item Correct Set ($S_{correct}$): Samples where the model's internal knowledge was sufficient ($\text{EM}(\hat{y}, y) = 1$). 
    \item Hallucination Set ($S_{hallucination}$): Samples where the model confidently generated an incorrect answer ($\text{EM}(\hat{y}, y) = 0$)
\end{itemize}

Using the partitioned data, we constructed a new supervised fine-tuning dataset (Dataset C) designed to reinforce known knowledge while penalizing hallucinations.

\textbf{Target Relabeling:} We applied a transformation function $T(x, \hat{y})$ to the targets based on the evaluation results:

$$T(x) =
\begin{cases}
\hat{y} & \text{if } x \in S_{correct} \quad \text{(Reinforcement)} \\
\text{"I don't know."} & \text{if } x \in S_{hallucination} \quad \text{(Abstention)}
\end{cases}$$

\note{Note: For the reinforcement targets, we utilized the model's own generated answers (Self-Training) to ensure the training distribution matched the model's preferred output style.}

\textbf{Class Balancing (80/20 Strategy)}: A significant risk in abstention training is the "laziness" objective, where a model learns that outputting "I don't know" is a global minimum for loss reduction. To prevent the model from defaulting to refusal, we enforced a strict class balance. We sampled a final dataset of 10,000 examples with a fixed ratio:
\begin{itemize}
    \item 80\% Reasoning Samples: Queries the model previously answered correctly.
    \item 20\% Refusal Samples: Queries where the model was taught to abstain.
\end{itemize}

This 4:1 ratio ensures the model prioritizes answering when confident, treating abstention as a fallback mechanism rather than a primary strategy. We then fine-tuned the Qwen2.5-7B-Instruct model using QLoRA (Quantized Low-Rank Adaptation) to maintain computational efficiency

-------------------------------------------------------------------------------\newline
\textbf{Reinforcement Learning from Verifier Feedback (RLVF):} To further improve the calibration of our abstention mechanism, we applied policy gradient reinforcement learning on top of the supervised abstention model (Qwen2.5-Abstain). While supervised fine-tuning teaches the model \textit{when} to abstain based on static training labels, we hypothesized that reinforcement learning could dynamically optimize the trade-off between attempting answers and abstaining based on reward signals tied to both correctness and confidence.

\textbf{Reward Architecture:} We designed a multi-component reward function $R(x, \hat{y}, y)$ that balances three objectives:

\begin{equation}
R(x, \hat{y}, y) = \begin{cases}
+10.0 & \text{if } \text{EM}(\hat{y}, y) = 1 \quad \text{(Correct Answer)} \\
-5.0 & \text{if } \hat{y} \neq y \text{ and } \hat{y} \neq \text{``I don't know''} \quad \text{(Hallucination)} \\
-1.0 & \text{if } \hat{y} = \text{``I don't know''} \quad \text{(Abstention Cost)}
\end{cases}
\end{equation}

This reward structure creates a natural hierarchy: correct answers are strongly rewarded ($+10$), hallucinations are heavily penalized ($-5$), and abstentions incur a small penalty ($-1$) to prevent the model from becoming overly conservative. The 10:5:1 ratio was calibrated to maintain the model's willingness to attempt questions while discouraging confident errors.

\textbf{Verifier-Based Factuality Check:} To detect hallucinations during training, we employed an NLI-based factuality verifier adapted from \cite{factualityverifier}. For each generated answer $\hat{y}$, we compute:
\begin{itemize}
    \item $p_{\text{entail}}$: Probability that $\hat{y}$ is entailed by the question context
    \item $p_{\text{contradict}}$: Probability that $\hat{y}$ contradicts known information
    \item Factuality score: $f = p_{\text{entail}} - p_{\text{contradict}}$
\end{itemize}

Answers with low factuality scores ($f < \theta$) trigger the hallucination penalty even if they superficially match formatting expectations. This prevents the model from learning to confidently output plausible-sounding but incorrect answers.
\textbf{Key Design Choices:} 
\begin{enumerate}
    \item \textbf{Warm-start from Supervised Model:} Rather than training from the base model, we initialized RLVF from Qwen2.5-Abstain. This allowed the model to start with a reasonable abstention boundary and focus RL on refining the decision threshold.
    \item \textbf{Asymmetric Penalties:} The 2× difference between hallucination penalty ($-5$) and abstention cost ($-1$) reflects our priority: it is better to abstain than to hallucinate, but the model should not default to constant abstention.
    \item \textbf{NLI-Based Verification:} Unlike reward models trained on human preferences, our verifier provides a deterministic, interpretable signal. This reduces the risk of reward hacking and ensures the model optimizes for factual correctness rather than stylistic preferences.
\end{enumerate}
---------------------------------------------------------------------------
\subsection{Scalable Implementation and Model Transition}

\paragraph{Platform Migration and Architecture Upgrade}
Initial experiments utilized Qwen2.5-7B on a standard virtual machine infrastructure. However, to address training instability and accelerate throughput (achieving a $\sim10\times$ speedup), we migrated our training pipeline to the \textbf{Tinker} high-performance compute platform. As the platform does not offer a Qwen2.5 model, we transitioned to \textbf{Qwen3-8B}. Validations confirmed that Qwen3-8B performs comparably to Qwen2.5-7B on our base benchmarks, ensuring that our architectural conclusions remain robust.

\paragraph{Data Scaling}
Leveraging this improved throughput, we scaled our supervised fine-tuning (SFT) from the initial 10,000-sample subset to the \textbf{full 80,000-sample HotPotQA training set}. Similarly, we re-distilled the CoT and Abstention models using the complete teacher-generated dataset, significantly expanding the breadth of training data compared to the pilot experiments.

\paragraph{Uncertainty Estimation: RAUQ $\to$ Semantic Entropy}
While our initial design proposed using the RAUQ metric \cite{vazhentsev2025rauq}, which relies on attention weight introspection, the abstracted inference API of our scaled platform precludes access to internal hidden states. Adopting a black-box approach, we implemented \textbf{Discrete Semantic Entropy (SE)} \cite{kuhn2023semantic}.

To estimate SE without access to token probabilities, we sample $M=10$ stochastic generations per question and cluster them based on semantic equivalence using a bidirectional NLI entailment model (\texttt{deberta-large-mnli}). The probability of each semantic cluster $C$ is approximated via Monte Carlo integration:
\begin{equation}
P(C|x) \approx \frac{1}{M} \sum_{i=1}^M \mathbb{I}(s_i \in C)
\end{equation}
The total uncertainty is then given by the entropy over these semantic clusters: $H(x) = -\sum_{C} P(C|x) \log P(C|x)$. This metric serves as our primary signal for evaluating the efficacy of our abstention mechanism in calculating and rejecting hallucinations.

\section{Experiments}
This section contains the following.

\subsection{Data}
Describe the dataset(s) you are using (provide references). If it's not already clear, make sure the associated task is clearly described.
Being precise about the exact form of the input and output can be very useful for readers attempting to understand your work, especially if you've defined your own task.

We use the the training and validation splits of the HotPotQA \cite{yang2018hotpotqadatasetdiverseexplainable} distractor dataset. The distractor setting has question-answer pairs with ten paragraphs as context, two of which are relevant to answering the question. Since we ran our tests in a closed-book setting, we discarded all ten context paragraphs for our finetuning, RL, and evaluation methods.

We define 3 training datasets:

\textbf{Dataset A: } The first 10,000 question-answer pairs of HotPotQA (distractor setting, train split).

\textbf{Dataset B:} The first 10,000 question-answer pairs of HotPotQA (distractor/train) with reasoning traces generated by the teacher model Qwen3-235B-A22B-Instruct-2507. The teacher model was given the two relevant paragraphs as context to reduce hallucinations in the generated reasoning traces. We prompted, filtered, and verified reasoning traces according to our "Finetune-CoT" section \ref{sec:finetune-cot}.
% Since this data is used to finetune the student model in a closed book setting, strict constraints were applied to the teacher's output to simulate "internalized knowledge." Specifically, the teacher model was prompted to avoid citation-style language (e.g., "according to the passage," "the text says") and instead frame the reasoning as a step-by-step recall of facts from memory. A robust validation loop was implemented to automatically reject and regenerate any traces that contained forbidden source-referencing phrases or failed to conclude with the exact ground truth answer.
\begin{itemize}
    \item \textbf{Q: } Which magazine was started first Arthur's Magazine or First for Women?
    \item \textbf{R: } 
    \begin{itemize} 
        \item Step 1: The question asks which magazine was started first between Arthur's Magazine and First for Women.
        \item Step 2: Arthur's Magazine was published starting in 1844, while First for Women, a modern woman's magazine, was launched much later, in the 1980s.
        \item Step 3: Since 1844 is earlier than the 1980s, Arthur's Magazine was started first. Therefore, the answer is Arthur's Magazine.
        \end{itemize}
    \item \textbf{A: } Arthur's Magazine


\end{itemize}

\textbf{Dataset C:} 10,000 question-answer or question-abstion pairs in the first 40,000 samples of HotPotQA (distractor/train). See our "Abstention-Aware Fine-Tuning" section \ref{sec:abstension} for details on generation.
\begin{itemize}
    \item \textbf{Q: } Who wrote the 1971 drama which stars the winner of the 1997 BAFTA Fellowship?
    
    \textbf{A: } I don't know.
    
    \item \textbf{Q: } Do Joshua Dugdale and Carlos Saldanha both work with films?
    
    \textbf{A: }yes
\end{itemize}

\subsection{Evaluation method}
Describe the evaluation metric(s) you use, plus any other details necessary to understand your evaluation.
Some projects will have clear metrics from prior work on given datasets, but we realize that other projects will define their own metrics.
If you're defining your own metrics, be clear as to what you're hoping to measure with each evaluation method (whether quantitative or qualitative, automatic or human-defined!), and how it's defined.

We use Exact Match and F1 between gold answers and the model's predicted answers to gauge performance. For abstention models, we also calculate Selective EM and Selective F1 by considering only the answers not abstained on. We also gauge the uncertainty of our models via RAUQ and semantic entropy during inference. By sorting model responses by their RAUQ score \cite{vazhentsev2025rauq} and iteratively rejecting the most uncertain samples, we use this as a target performance for our model's abstention.

calculated RAUQ, 

\subsection{Experimental details}
Report how you ran your experiments (e.g. model configurations, learning rate, training time, etc.)

\textbf{Model Configuration:} 
\begin{itemize}
    \item Qwen2.5-Instruct: Evaluated Qwen2.5-Instruct in its original, pre-trained state without any additional fine-tuning. The model was loaded using bfloat16 precision to optimize memory usage while maintaining numerical stability.
    \item Qwen2.5-SFT \& Qwen2.5-FCoT \& Qwen2.5-Abstain: To maintain computational efficiency, we utilized QLoRA (Quantized Low-Rank Adaptation). The base model was loaded in 4-bit NormalFloat (NF4) precision with double quantization, while the LoRA adapters were trained in FP16.
\end{itemize}

\textbf{Hyperparameters \& Training Configuration}
\begin{itemize}
    \item Qwen2.5-SFT \& Qwen2.5-FCoT \& Qwen2.5-Abstain: 
        \begin{itemize}
            \item LoRA Config: We applied LoRA adapters to all linear projection layers $(q_{proj}, k_{proj}, v_{proj}, o_{proj}, gate_{proj}, up_{proj}, down_{proj})$ with a rank $r=16$, alpha $\alpha=32$, and a dropout rate of $0.05$. 
            \item Optimization: The model was trained using the adamw\_torch optimizer with a learning rate of 2e-4. We used a cosine learning rate scheduler with a 3\% warmup ratio.
            \item Batching: Training was performed with a per-device batch size of 4 and gradient accumulation of 4 steps, resulting in an effective batch size of 16.
        \end{itemize}

    \item \textbf{Qwen2.5-SFT:}
        \begin{itemize}
            \item Formatting: Standard Q\&A pairs using the chat template.
            \item Max Length: 512 tokens.
            \item Masking: System prompt and user question were masked; loss was calculated only on the answer.
        \end{itemize}
    
    \item \textbf{Qwen2.5-FCoT \& Qwen2.5-Abstain:}
        
        \begin{itemize}
            \item Formatting: We adopted the delimiter style from "Language Models are Reasoning Teachers" to minimize token overhead. 
            \item Structure: \texttt{Question \#\#\# Rationale --> Answer END}.
            \item Max Length: Increased to 2048 tokens to accommodate full reasoning traces.
            \item Masking: The prompt (\texttt{Question \#\#\#}) was masked. The model was trained to generate both the reasoning trace and the final answer.
        \end{itemize}
    \item \textbf{REINFORCE Qwen2.5-Abstain}
    
    \begin{itemize}
        \item \textbf{Policy Initialization:} Started from the supervised abstention checkpoint to leverage its learned abstention boundary.
        \item \textbf{Optimization:} Used AdamW with learning rate $\eta = 1\times10^{-5}$ and weight decay $0.01$.
        \item \textbf{LoRA Configuration:} Applied rank-16 LoRA adapters to all linear layers with $\alpha=32$, maintaining the same adapter structure as the supervised models.
        \item \textbf{Training Data:} Sampled from the HotPotQA training split (distractor setting), using on-policy generation to maintain distribution alignment.
        \item \textbf{Training Duration:} 1 epoch over 9,992 gradient steps with batch size 4 and gradient accumulation of 4 steps (effective batch size 16).
    \end{itemize}
    

        

\end{itemize}

\textbf{Inference Setup: } All experiments were conducted in a strictly closed-book setting (no supporting paragraphs provided). 
\begin{itemize}
    \item Qwen2.5-Instruct and Qwen2.5-SFT were provided with the following system prompt:
        \begin{quote}
            You are a concise encyclopedia. Answer the question directly with a short phrase or entity name. Do not explain.
        \end{quote}
    This enforced brevity.
    
    \item Qwen2.5-FCoT \& Qwen2.5-Abstain: Prompted with the question followed by the delimiter \texttt{ \#\#\#}. The model generated a reasoning trace followed by \texttt{-->} and the answer. We programmatically parsed the final answer by splitting the output at the \texttt{-->} token
\end{itemize}

\textbf{Inference Configuration:} During evaluation, we used greedy decoding to ensure deterministic outputs. 
\begin{itemize} 
    \item For standard models (Instruct/SFT/Abstain), the maximum generation length was 50 tokens. \item For the CoT model, the maximum generation length was increased to 512 tokens to allow for reasoning generation. 
\end{itemize} 
To enable the RAUQ calculation, we forced the model to output full attention weights and used the eager attention implementation. Inference was run with a batch size of 4 to accommodate the significant VRAM overhead incurred by storing full attention matrices.





\subsection{Results}
Report the quantitative results that you have found so far. Use a table or plot to compare results and compare against baselines.

\begin{itemize}
    \item Comment on your quantitative results. Are they what you expected? Better than you expected? Worse than you expected? Why do you think that is? What does that tell you about your approach?
\end{itemize}

\textbf{Generation Performance and Error Correction:} Our fine-tuning experiments demonstrate clear improvements over the Qwen2.5-Instruct baseline across all metrics. As shown in Table \ref{tab:model_performance}, the Chain-of-Thought (CoT) model achieved the highest exact match accuracy of 24.31\%, a significant gain over the Base model's 18.86\%. Notably, the CoT approach proved superior in error correction, successfully rectifying 697 questions that the Base model originally answered incorrectly, compared to 477 corrections by the SFT model. While both fine-tuning methods introduced regressions—breaking a smaller number of previously correct answers (249 for SFT and 294 for CoT)—the net performance gain validates the efficacy of distilling reasoning traces over standard supervised fine-tuning for this domain.

\textbf{Abstention Mechanism and Reliability:} To mitigate hallucinations, we developed two complementary abstention models using different training paradigms. We first trained a supervised abstention model (Qwen2.5-Abstain) by relabeling questions where the CoT baseline failed with "I don't know" targets, using an 80/20 balance to prevent over-abstention. We then applied policy gradient reinforcement learning (Qwen2.5-RLVF) on top of this supervised checkpoint to refine the abstention boundary using reward shaping.

Table \ref{tab:abstention_analysis} presents their precision and recall metrics, where Precision measures the percentage of abstentions that correspond to questions the reference model answered incorrectly (validating that the model abstains when appropriate), and Recall measures the percentage of reference model errors that were successfully avoided through abstention.

The supervised model (Qwen2.5-Abstain) demonstrates highly conservative behavior, achieving 97.49\% precision against the Base model, meaning its abstentions are almost invariably justified. Its recall ranges from 55-61\% across reference models, performing best on the "Hardest" subset (60.85\%)—questions where all baseline models failed. This confirms the mechanism activates most strongly on genuinely difficult queries.

The RLVF model exhibits a different profile: it exhibitds higher precision (91-98\%) but achieves lower recall (50-54\%). This trade-off is intentional—the reward structure penalizes abstention modestly ($-1$) relative to hallucinations ($-5$), encouraging the model to attempt more questions. As shown in Table \ref{tab:hotpotqa_comparison}, this yields higher coverage: the RLVF model answers 58.83\% of questions versus the supervised model's 53.65\%, while maintaining identical overall EM (20.37\%).

However, this increased coverage comes at the cost of selective accuracy. The supervised model achieves 37.98\% EM on attempted questions, compared to the RLVF model's 34.62\%. This reflects the fundamental trade-off visualized in Figure~\ref{fig:rlhf_confusion_matrices}: the supervised approach prioritizes safety (high precision, conservative abstention), while the RLVF approach prioritizes coverage (more attempts, lower selective accuracy). Both strategies are valid depending on application requirements—safety-critical domains benefit from the supervised model's conservatism, while general-purpose systems benefit from the RLVF model's higher answer rate.

% \begin{table}[h]
% \centering
% \caption{Head-to-Head Comparison: Qwen2.5-SFT vs. Qwen2.5-FCoT}
% \label{tab:sft_vs_cot}
% \begin{tabular}{lr}
% \toprule
% \textbf{Comparison Case} & \textbf{Count} \\
% \midrule
% Both Correct & 1,244 \\
% Both Wrong & 5,224 \\
% CoT Corrected (SFT Wrong $\to$ CoT Right) & 556 \\
% SFT Prevailed (SFT Right $\to$ CoT Wrong) & 381 \\
% \bottomrule
% \end{tabular}   
% \end{table}
\begin{table}[H]
\centering
\caption{Comparison of Experimental and State-of-the-Art Results on the HotpotQA Dataset (Distractor/Validation) (EM and F1 scores are percentages).}
\label{tab:hotpotqa_comparison}
\begin{tabular}{lrrrrrr}
    \toprule
    \textbf{Model} & \textbf{EM} & \textbf{F1} & \textbf{Abs. Rate} & \textbf{Sel. EM} & \textbf{Sel. F1} & \textbf{Avg. RAUQ} \\
    \midrule
    \multicolumn{7}{l}{\textbf{Experimental Results ($N=7405$ Samples)}} \\
    \cmidrule(lr){1-7}
    Qwen2.5-Instruct & 18.86 & 27.02 & -- & -- & -- & 1.25 \\
    Qwen2.5-SFT & 21.94 & 31.41 & -- & -- & -- & 1.62 \\
    Qwen2.5-FCoT & \textbf{24.31} & 33.87 & -- & -- & -- & -- \\
    Qwen2.5-Abstain & 20.37 & 26.04 & 46.34 & 37.98 & 48.46 & 1.32 \\
    Qwen2.5-RLVF & 20.37 & 26.24 & 41.17 & 34.62 & 44.60 & 1.40 \\
    \midrule
    \multicolumn{7}{l}{\textbf{SOTA Closed-Book Results (RECITE Paper)}} \\
    \cmidrule(lr){1-7}
    Codex-002 & 37.11 & 48.37 & N/A & N/A & N/A & N/A \\
    PaLM-62B & 26.46 & 35.67 & N/A & N/A & N/A & N/A \\
    \bottomrule
\end{tabular}
\end{table}
\begin{figure}[H]
\centering

% SFT Abstention Model (Top)
\includegraphics[width=1\textwidth]{abstention_confusion_matrices_SFT.png}
\vspace{0.5cm}

% RLHF Model (Bottom)
\includegraphics[width=1\textwidth]{rlhf_confusion_matrices.png}

\caption{Abstention Performance Comparison: Supervised Fine-Tuning vs. RLHF. \textbf{Top}: Confusion matrices for the supervised abstention model (Qwen2.5-Abstain SFT), which was trained with explicit abstention labels on questions where baseline models failed. \textbf{Bottom}: Confusion matrices for the RLHF model, trained using policy gradient reinforcement learning with reward shaping. Each panel shows how well the model's abstention decisions align with questions that reference models answered incorrectly. }
\end{figure}

\label{fig:rlhf_confusion_matrices}
\end{figure}
\begin{table}[H]
\centering
\caption{Performance Analysis of Abstention Mechanism Across Reference Models}
\label{tab:abstention_analysis}
\resizebox{\columnwidth}{!}{%
\begin{tabular}{lcccccc}
\toprule
\textbf{Model} & \textbf{Ref. Model} & \textbf{Precision} & \textbf{Recall} & \textbf{TP} & \textbf{FP} & \textbf{FN} \\
 & & & & (Good Abs) & (Bad Abs) & (Missed) \\
\midrule
\multirow{}{Qwen2.5-Abstain} 
    & Base & \textbf{97.49\%} & 55.69\% & 3,346 & \textbf{86} & 2,662 \\
    & SFT & 95.72\% & 56.83\% & 3,285 & 147 & 2,495 \\
    & CoT & 93.97\% & 57.54\% & 3,225 & 207 & 2,380 \\
    & Hardest & 90.53\% & \textbf{60.85\%} & 3,107 & 325 & \textbf{1,999} \\
\midrule
\multirow{Qwen2.5-RLVF}
    & Base & \textbf{97.84\%} & 49.65\% & 2,983 & \textbf{66} & 3,025 \\
    & SFT & 96.23\% & 50.76\% & 2,934 & 115 & 2,846 \\
    & CoT & 94.19\% & 51.24\% & 2,872 & 177 & 2,733 \\
    & Hardest & 91.21\% & 54.47\% & 2,781 & 268 & 2,325 \\
\bottomrule
\end{tabular}%
}
\end{table}




% \begin{table*}[h]
% \centering
% \caption{Combined Performance Analysis: Generation Accuracy vs. Abstention Efficacy}
% \label{tab:combined_performance}
% \resizebox{\textwidth}{!}{% Ensures table fits within text width
% \begin{tabular}{lccccc}
% \toprule
%  & \multicolumn{3}{c}{\textbf{Generation Performance}} & \multicolumn{2}{c}{\textbf{Abstention Performance}} \\
% \cmidrule(lr){2-4} \cmidrule(lr){5-6}
% \textbf{Model / Setting} & \textbf{Accuracy} & \textbf{Corrected} & \textbf{Regressed} & \textbf{Precision} & \textbf{Recall} \\
%  & (Exact Match) & (vs Base) & (vs Base) & & \\
% \midrule
% \textbf{Base Model} & 18.87\% & -- & -- & \textbf{97.49\%} & 55.69\% \\
% \textbf{SFT Model} & 21.94\% & 477 & \textbf{249} & 95.72\% & 56.83\% \\
% \textbf{CoT Model} & \textbf{24.31\%} & \textbf{697} & 294 & 93.97\% & 57.54\% \\
% \textbf{Hardest Subset} & -- & -- & -- & 90.53\% & \textbf{60.85\%} \\
% \bottomrule
% \end{tabular}%
% }
% \end{table*}

\begin{table}[h]
\centering
\caption{Overall Model Performance and Error Analysis Relative to Base Model}
\label{tab:model_performance}
\begin{tabular}{lccc}
\toprule
\textbf{Metric} & \textbf{Base Model} & \textbf{SFT Model} & \textbf{CoT Model} \\
\midrule
Overall Accuracy (Exact Match) & 18.87\% & 21.94\% & \textbf{24.31\%} \\
Corrections (Base Wrong $\to$ Right) & -- & 477 & \textbf{697} \\
Regressions (Base Right $\to$ Wrong) & -- & \textbf{249} & 294 \\
\bottomrule
\end{tabular}
\end{table}

\subsubsection{Semantic Entropy and Abstention Results}

% Row 1: Distributions (Box/Violin Plots)
\begin{figure}[H]
    \centering
    \setlength{\abovecaptionskip}{2pt} % Reduce caption padding
    \begin{minipage}{0.32\textwidth}
        \centering
        \includegraphics[height=3.5cm, keepaspectratio]{Semantic Entropy/SE ABS fixed SE (400)/se_correlation_boxplot_std.png}
        \caption*{Std SE: Abstained vs Answered}
    \end{minipage}%
    \hfill
    \begin{minipage}{0.32\textwidth}
        \centering
        \includegraphics[height=3.5cm, keepaspectratio]{Semantic Entropy/SE ABS fixed SE (400)/se_correlation_boxplot_cond.png}
        \caption*{Cond SE: Abstained vs Answered}
    \end{minipage}%
    \hfill
    \begin{minipage}{0.32\textwidth}
        \centering
        \includegraphics[height=3.5cm, keepaspectratio]{Semantic Entropy/SE ABS fixed SE (400)/se_correctness_violin.png}
        \caption*{Std SE: Correct vs Incorrect}
    \end{minipage}
    \caption{Distribution of Uncertainty (Entropy) across outcomes.}
    \label{fig:se_distributions}
    \vspace{-10pt} % Pull next content up
\end{figure}

\begin{figure}[H]
    \centering
    \setlength{\abovecaptionskip}{2pt}
    \begin{minipage}{0.48\textwidth}
        \centering
        \includegraphics[height=5cm, keepaspectratio]{Semantic Entropy/SE ABS fixed SE (400)/se_curves_full_f1_score.png}
        \caption*{Full Population (Efficiency)}
    \end{minipage}%
    \hfill
    \begin{minipage}{0.48\textwidth}
        \centering
        \includegraphics[height=5cm, keepaspectratio]{Semantic Entropy/SE ABS fixed SE (400)/se_curves_selective_f1_score.png}
        \caption*{Selective Population (Risk Mitigation)}
    \end{minipage}
    \caption{F1 Score Rejection Curves: Comparing efficiency on the full dataset vs risk mitigation on answered questions.}
    \label{fig:se_f1_curves}
    \vspace{-10pt}
\end{figure}

\paragraph{Uncertainty Quantification vs. Correctness}
We validated the efficacy of Semantic Entropy (SE) as a proxy for correctness. As shown in Figure \ref{fig:se_distributions} (Right), the distribution of SE for correct answers is heavily clustered around 0 (median $\approx 0$), whereas incorrect answers exhibit significantly higher variance with a median of $\sim 0.9$. This confirms that high Semantic Entropy successfully signals likely hallucinations.

\paragraph{Abstention Calibration}
We analyzed whether the model's native abstention mechanism aligns with semantic uncertainty (Figure \ref{fig:se_distributions} Left/Center). For Standard SE, the ``Abstained'' distribution (Median $\approx 0.65$) is consistently shifted higher than the ``Answered'' distribution (Median $\approx 0.5$). This indicates the model is correctly choosing to abstain on questions where it is internally uncertain, though the overlap suggests further gains are possible by explicitly rejecting high-entropy answers.

\paragraph{Rejection Curve Efficiency}
Figure \ref{fig:se_f1_curves} demonstrates the trade-off efficiency of utilizing SE for rejection. The \textbf{Consensus} strategy (Sampled) consistently outperforms the \textbf{Greedy} strategy across all rejection thresholds. 
\begin{itemize}
    \item \textbf{Base Performance:} At 0\% rejection, Consensus achieves a Selective F1 of 0.50 compared to Greedy's 0.39 (+28\% relative improvement).
    \item \textbf{Rejection Gain:} By rejecting the top 50\% most uncertain samples, the Consensus Selective F1 score improves to \textbf{0.69}, a substantial gain in reliability for safety-critical applications.
\end{itemize}

\section{Analysis}

\subsection{"Reasoning Trace" Qualitative Check}
\textbf{Evidence of Successful "Internalized Knowledge" Transfer:} Qualitative inspection reveals that the CoT model successfully learned the 'internalized knowledge' style, effectively using the reasoning steps to bridge gaps where the Base model failed.

In the example regarding the shared nationality of Scott Derrickson and Ed Wood (ID: 5a8b57f2...), the Base model immediately predicted the incorrect answer No. However, the CoT model generated a robust 3-step trace:

Question: Were Scott Derrickson and Ed Wood of the same nationality?
Gold Answer: yes
Base Pred (Wrong): No
CoT Pred (Correct): yes
-------------------- Reasoning Trace --------------------
Step 1: The question asks about the nationalities of Scott Derrickson and Ed Wood.  
Step 2: Scott Derrickson is an American film director, born in the United States, while Ed Wood was an American filmmaker known for his work in B-movies.  
Step 3: Since both individuals were born and active in the United States, they share the same nationality. Therefore, the answer is yes.

Notably, the student model adhered to the negative constraint, avoiding citations or phrases like "According to the text," suggesting successful stylistic transfer from the teacher while correcting the factual error.

\textbf{Analysis of Regression \& Hallucination} However, the reasoning process introduces two distinct failure modes where the Base model was originally correct but the CoT model regressed.

First, the step-by-step generation can occasionally induce hallucinations that derail the answer. For instance, in the query about Rose McGowan, the Base model correctly answered Charmed via simple association. The CoT model, however, attempted to reason through her roles and hallucinated a specific character:

Question: What WB supernatrual drama series was Jawbreaker star Rose Mcgowan best known for being in?
Gold Answer: Charmed
Base Pred (Correct): Charmed
CoT Pred (Wrong): Buffy the Vampire Slayer
-------------------- Reasoning Trace --------------------
Step 1: The question asks about the WB supernatural drama series that Rose McGowan is best known for starring in.  
Step 2: Rose McGowan is best known for her role as Cordelia Chase in the WB supernatural drama series "Buffy the Vampire Slayer."  
Therefore, the answer is Buffy the Vampire Slayer.


Second, the CoT model sometimes suffers from 'granularity mismatches' where the reasoning leads to higher specificity than the gold label. When asked about the occupation of Chris Menges and Aram Avakian, the Base model correctly predicted Director. The CoT model, identifying their specific bodies of work, deduced:

Question: What occupation do Chris Menges and Aram Avakian share?
Gold Answer: director
Base Pred (Correct): Director
CoT Pred (Wrong): film director
-------------------- Reasoning Trace --------------------
Step 1: The question asks about the shared occupation of Chris Menges and Aram Avakian.  
Step 2: Chris Menges is known as a British film director, and Aram Avakian is an American film director and producer.  
Step 3: Since both individuals are primarily recognized for directing films, their shared occupation is film director. Therefore, the answer is film director.


While factually correct, the prediction film director was penalized against the rigid gold label director, demonstrating that the model’s detailed reasoning can lead to over-specification that strict Exact Match (EM) metrics fail to capture, hence the benefit of including F1 score in our evaluation metrics.

\subsection{SFT vs.\ RLVF Abstention Behavior (Risk Comparison)}

Comparing the two abstention strategies reveals clear behavioral boundaries. For the query, \emph{``What government position was held by the woman who portrayed Corliss Archer in the film \textit{Kiss and Tell}?''}, the SFT model responded with ``I don't know,'' reflecting the conservative abstention style learned during supervised training. In contrast, the RLVF model, which is optimized under an asymmetric reward structure that penalizes abstention ($-1$) less severely than incorrect answers ($-5$), chose to attempt an answer and predicted ``First Lady.'' Because the correct answer is ``Chief of Protocol,'' this instance illustrates a \textit{confident hallucination}: the model preferred risking an incorrect response rather than accepting the abstention penalty. Such cases demonstrate the ``risk'' side of the trade-off, where the model's incentive to avoid the abstention penalty leads to unreliable or overconfident guesses.

However, the same mechanism can also produce beneficial behavior. For example, in the query \emph{``What army did the namesake of the ship launched as the M\"unchen in 1930 fight in during the American Revolutionary War?''}, the SFT model again abstained with ``I don't know.'' The RLVF model, facing the same penalty structure, elected to answer and produced the correct response: ``Continental Army.'' This represents the ``reward'' side of the trade-off: the model takes a calculated risk and, in this case, succeeds in recovering information that the SFT baseline refuses to commit to. Instances like this show that the asymmetric reward can indeed encourage productive non-abstentions, enabling the model to fill gaps where supervised training remains overly cautious.

Together, these two examples highlight the central tension of abstention-aware RLHF. A more aggressive penalty on abstention encourages the model to answer more often, which can yield legitimate gains in coverage, but also increases the likelihood of confident hallucinations. Achieving the right balance requires careful tuning of the abstention penalty and, potentially, more structured confidence modeling to distinguish informative risk-taking from unreliable guessing.


\section{Conclusion}
\section{Conclusion}
In this work, we presented a comprehensive approach to mitigating hallucination in Small Language Models for closed-book question answering. By decomposing the hallucination problem into "access failures" and "knowledge gaps," we applied targeted interventions: Chain-of-Thought distillation to improve reasoning-based retrieval, and Abstention-Aware Fine-Tuning to calibrate the model's refusal behavior.

Our experiments on HotPotQA reveal that these strategies are highly effective. CoT distillation proved superior to standard fine-tuning, correcting 697 baseline errors by guiding the model to step-by-step derivational answers. For uncorrectable knowledge gaps, our supervised abstention mechanism achieved a 97.5\% precision in flagging hallucinations, effectively converting "confident errors" into "safe refusals." Furthermore, we demonstrated that Reinforcement Learning with Verifier Feedback (RLVF) offers a dynamic method to tune the trade-off between caution and helpfulness, allowing system designers to optimize for coverage without sacrificing fundamental reliability.

However, limitations remain. Our RLVF approach, while increasing coverage, slightly degraded selective accuracy compared to the purely supervised baseline, highlighting the difficulty of optimizing reward landscapes where "safety" (abstention) is a local optimum. Additionally, our reliance on Semantic Entropy for post-hoc filtering, while effective, incurs a significant computational cost during inference.

Future work will focus on two directions: (1) developing lightweight, cascading uncertainty filters (e.g., using RAUQ as a pre-filter for Semantic Entropy) to reduce inference latency, and (2) extending our evaluation to open-book settings to assess whether these calibration gains transfer when the model must discriminate between parametric knowledge and retrieved context. Ultimately, our findings suggest that "knowing what you don't know" is a learnable capability for SLMs, paving the way for more trustworthy autonomous agents.

\section{Team Contributions}
\section{Team Contributions}
\begin{itemize}
    \item \textbf{Joshua Hegstad:} Led the implementation of the Reinforcement Learning (RLVF) pipeline and the Semantic Entropy evaluation framework. Conducted the primary analysis of risk-coverage trade-offs and authored the report's result analysis.
    \item \textbf{Ahmed Jaber:} Developed the Chain-of-Thought (CoT) distillation pipeline, including teacher prompting, trace filtering, and student fine-tuning. Managed the transition to the Tinker high-performance compute platform.
    \item \textbf{Farhaan Siddiqui:} Curated the HotPotQA datasets and implemented the supervised fine-tuning (SFT) baselines. Performed the initial RAUQ uncertainty integration and managed data partitioning to prevent leakage.
\end{itemize}


\bibliographystyle{unsrt}
\bibliography{references}

\appendix

\section{Appendix (optional)}
If you wish, you can include an appendix, which should be part of the main PDF, and does not count towards the 6-8 page limit.
Appendices can be useful to supply extra details, examples, figures, results, visualizations, etc., that you couldn't fit into the main paper. However, your grader \textit{does not} have to read your appendix, and you should assume that you will be graded based on the content of the main part of your paper only.

\subsection{Additional Rejection Curves}

\begin{figure}[H]
    \centering
    \setlength{\abovecaptionskip}{2pt}
    \begin{minipage}{0.48\textwidth}
        \centering
        \includegraphics[height=5cm, keepaspectratio]{Semantic Entropy/SE ABS fixed SE (400)/se_curves_full_exact_match.png}
        \caption*{Full Population (Efficiency)}
    \end{minipage}%
    \hfill
    \begin{minipage}{0.48\textwidth}
        \centering
        \includegraphics[height=5cm, keepaspectratio]{Semantic Entropy/SE ABS fixed SE (400)/se_curves_selective_exact_match.png}
        \caption*{Selective Population (Risk Mitigation)}
    \end{minipage}
    \caption{Exact Match (EM) Rejection Curves.}
    \label{fig:se_em_curves}
    \vspace{-10pt}
\end{figure}

\begin{table}[H]
\centering
\caption{Impact of Semantic Entropy-based Rejection on F1 Scores ($N=400$)}
\label{tab:se_rejection_stats}
\begin{tabular}{llrr}
\toprule
\textbf{Strategy} & \textbf{Rejection Threshold} & \textbf{Full F1} & \textbf{Selective F1} \\
\midrule
Greedy     & 0\%          & 0.2200 & 0.3894 \\
Greedy     & 50\%         & 0.3228 & 0.5713 \\
\midrule
Consensus  & 0\%          & 0.2645 & 0.4991 \\
Consensus  & 50\%         & 0.3725 & \textbf{0.6872} \\
\bottomrule
\end{tabular}
\end{table}

\end{document}